\chapter{Hernia Repair}
\index{Hernia Repair}

\section*{Definition and Overview}
Hernia repair is surgery to correct a hernia, in which an organ or tissue protrudes through a weak spot in the muscle or connective tissue that normally holds it in place. The most common hernias are inguinal (groin), umbilical (belly button), incisional (through a previous scar), and hiatal (stomach through the diaphragm). Repair involves returning the protruding tissue to its proper position and reinforcing the weakened area, usually with a synthetic mesh. Most repairs are performed as day surgery using keyhole (laparoscopic) or open techniques, with excellent results and a quick return to normal activity.

\section*{Epidemiology}
Hernias are very common: inguinal hernias affect around 25--40\% of men over their lifetime and are the most frequently repaired condition in general surgery worldwide. Approximately 80,000 hernia repairs are performed in many countries each year. Inguinal hernias are far more common in men, while femoral hernias are more common in women and carry a higher strangulation risk. Umbilical hernias affect many newborns but usually close spontaneously, while those that persist or appear in adults require repair.

\section*{Aetiology and Pathophysiology}
\begin{itemize}
    \item \textbf{Weakness:} congenital defects, ageing, previous surgery, and connective tissue factors weaken the abdominal wall
    \item \textbf{Increased pressure:} lifting, straining, chronic cough, constipation, obesity, and pregnancy raise intra-abdominal pressure
    \item \textbf{Inguinal hernias:} protrusion through the inguinal canal, mainly in men
    \item \textbf{Femoral hernias:} through the femoral canal, mainly in women, with higher strangulation risk
    \item \textbf{Umbilical hernias:} through the belly button, common in babies and in adults with abdominal pressure
    \item \textbf{Incisional hernias:} through a weakened surgical scar
    \item \textbf{Hiatal hernias:} part of the stomach slides up through the diaphragm into the chest
    \item \textbf{Complications:} the hernia contents can become stuck (incarcerated) and lose blood supply (strangulated), a surgical emergency
\end{itemize}

\section*{Clinical Features}
\begin{itemize}
    \item A bulge that appears with straining, coughing, or standing and reduces when lying down
    \item Aching, discomfort, or a dragging sensation at the site
    \item Pain on lifting or exertion
    \item Inguinal hernia: a bulge in the groin, sometimes extending into the scrotum
    \item Umbilical hernia: a bulge at the belly button
    \item Incarceration: a lump that cannot be pushed back, with increasing pain
    \item Strangulation: severe pain, redness, nausea, vomiting, and fever, requiring emergency surgery
\end{itemize}

\section*{Differential Diagnosis}
\begin{itemize}
    \item Hydrocoele and testicular swellings in the groin
    \item Enlarged lymph nodes and cysts in the groin
    \item Varicocoele and other scrotal conditions
    \item Muscle strain and sports hernia (athletic pubalgia)
    \item Femoral artery aneurysm
    \item Hiatal hernia is diagnosed by endoscopy and imaging, presenting with reflux symptoms
    \item Clinical examination and ultrasound usually distinguish these
\end{itemize}

\section*{When to Seek Medical Care}
See a doctor for any new or changing lump, bulge, or discomfort in the groin, belly button, or an old scar. A painful hernia that cannot be reduced requires urgent assessment.

\textbf{Contact your local emergency services for:}
\begin{itemize}
    \item A hernia that becomes painful, tender, red, or cannot be pushed back
    \item Severe abdominal pain with nausea and vomiting (possible strangulation)
    \item Fever with a tender lump
    \item Vomiting with abdominal distension and inability to pass wind or stool
\end{itemize}

\section*{Diagnosis and Investigations}
\begin{itemize}
    \item \textbf{Clinical examination:} most hernias are diagnosed by examination, including feeling for a cough impulse
    \item \textbf{Ultrasound:} confirms the diagnosis when examination is inconclusive
    \item \textbf{CT scan:} used for complex, recurrent, or uncertain hernias
    \item \textbf{Endoscopy and barium studies:} for hiatal hernias
    \item \textbf{Preoperative assessment:} fitness for surgery, including anaesthetic review
    \item \textbf{Urgent assessment:} incarcerated or strangulated hernias need emergency imaging and surgery
\end{itemize}

\section*{Management}
\begin{itemize}
    \item \textbf{Observation:} small, symptom-free hernias in some adults may be safely watched
    \item \textbf{Open repair:} a direct incision over the hernia, with the defect closed and reinforced with mesh
    \item \textbf{Laparoscopic repair:} keyhole surgery through small incisions, with faster recovery and less pain
    \item \textbf{Mesh:} synthetic mesh is used in most repairs to reduce recurrence
    \item \textbf{Day surgery:} most repairs are performed as day procedures with local or general anaesthesia
    \item \textbf{Emergency repair:} incarcerated or strangulated hernias require urgent surgery, sometimes with bowel resection
    \item \textbf{Recovery:} return to light activity within days and heavy lifting within 4--6 weeks
\end{itemize}

\section*{Complications}
\begin{itemize}
    \item Strangulation, with gangrene of the bowel and life-threatening sepsis if untreated
    \item Recurrence, reduced with mesh repair but still occurring in a small percentage
    \item Chronic pain in the groin after inguinal repair
    \item Bruising, swelling, and wound infection
    \item Seroma or haematoma formation at the repair site
    \item Nerve injury and numbness
    \item Risks of anaesthesia and, rarely, injury to surrounding structures
\end{itemize}

\section*{Prognosis and Outcomes}
Hernia repair is one of the most successful operations in surgery. More than 95\% of repairs are complication-free, recurrence rates with mesh are under 5\%, and most people return to all normal activities within weeks. Laparoscopic repair offers less postoperative pain and faster recovery. The main long-term issues are chronic groin pain in a minority and the small recurrence risk. Repairing hernias before they strangulate avoids the serious complications of emergency surgery.

\section*{Prevention and Self-Care}
\begin{itemize}
    \item See a doctor promptly for any bulge so repair can be planned before complications
    \item Use proper lifting technique and avoid sudden heavy strain
    \item Prevent and treat constipation with fibre, fluids, and activity
    \item Treat chronic cough and manage obesity
    \item Follow the surgeon's recovery advice, including graduated lifting
    \item Report any recurrence or new bulge after surgery
    \item After repair, strengthen core muscles gradually under guidance
\end{itemize}

\section*{Sources and Further Reading}
The following reputable sources provide further information for healthcare students and patients seeking reliable, current advice about hernia repair.

\begin{itemize}
    \item Healthdirect. \textit{Hernia repair: what to expect.} https://www.healthdirect.gov.au/hernia-repair
    \item Royal Australasian College of Surgeons. \textit{Hernia repair information.} https://www.surgeons.org/
    \item Australian Hernia Society. \textit{Patient information about hernias.} https://www.hernia.org.au/
    \item NHS. \textit{Hernia repair: overview.} https://www.nhs.uk/conditions/hernia/
    \item Mayo Clinic. \textit{Hernia repair: why it's done.} https://www.mayoclinic.org/
    \item MedlinePlus. \textit{Hernia repair.} https://medlineplus.gov/hernia.html
\end{itemize}
